\section{A multiplicity-one family detected by a Hessian bracket}

\subsection*{Finding F-024: an exact near-rectangular plethysm coefficient}
For every \(d\geq3\) and \(n\geq3\), put
\[
                     \lambda=(d(n-2)+3,3,2^{d-3}).
                                                               \tag{F-024.1}
\]
Then, in characteristic zero,
\begin{equation}\label{eq:F024ambient}
             a_\lambda(d[n])=1.
\end{equation}
The ambient vector space must have dimension at least
\(d-1=\ell(\lambda)\); otherwise this Schur module is zero.
This exact extraction is derived below from standard Cauchy
and Pieri identities, without a claim of historical priority.
It concerns the complete ambient multiplicity, rather than
only the existence of one highest-weight vector.

\paragraph{The column representation and its balanced weight.}
Write \(a=n-2\geq1\), \(D=da\), and let
\(V=\mathbb C^d\) be the column space. The Cauchy
polynomial-matrix model used in F-014 and F-016 identifies
\(a_\lambda(d[n])\) with
\[
                 \dim\left((S_\lambda V)_{(n^d)}\right)^{S_d}.
                                                               \tag{F-024.2}
\]
Indeed the degree \((n,\ldots,n)\) column multigrading
is \((\operatorname{Sym}^n U)^{\otimes d}\), and
column-permutation invariants are
\(\operatorname{Sym}^d\operatorname{Sym}^n U\).
Cauchy identifies the complete row highest-weight space
with the column module \(S_\lambda V\), preserving that
multigrading and the column permutations.
The precise standard inputs are Macdonald, I(4.3), p.~63,
and Appendix~A(5.3)--(5.6), pp.~155--156.

Pad \(\lambda\) with a final zero and twist this column
module by \(\det(V)^{-2}\). Its highest weight becomes
\[
                   \theta=(D+1,1,0^{d-3},-2),
                                                               \tag{F-024.3}
\]
and the balanced weight becomes \((a^d)\).
Every column permutation has determinant \(\pm1\),
so its action on the determinant-square twist is trivial.
Therefore (F-024.2) is exactly
\(\dim(L_\theta{}_{(a^d)})^{S_d}\), where \(L_\theta\)
is the irreducible rational \(\mathrm{GL}(V)\)-module
of highest weight (F-024.3).

\paragraph{A four-term rational character identity.}
Write \(h_j\) for \(\operatorname{Sym}^j V\),
write \(V^*\) for the dual standard module, and interpret
the following identities in the characteristic-zero
representation Grothendieck group. Dual Pieri gives
\[
\begin{split}
 [S_{(D+1,1)}V\otimes\operatorname{Sym}^2V^*]
  ={}&[L_\theta]
       +[L_{(D,1,0^{d-3},-1)}]
       +[L_{(D+1,0^{d-2},-1)}]\\
     &+[S_{(D-1,1)}V]+[h_D],\\
 [S_{(D,1)}V\otimes V^*]
  ={}&[L_{(D,1,0^{d-3},-1)}]
       +[S_{(D-1,1)}V]+[h_D],\\
 [h_{D+1}\otimes V^*]
  ={}&[L_{(D+1,0^{d-2},-1)}]+[h_D].
\end{split}\tag{F-024.4}
\]
For clarity, these are the complete dual Pieri lists, not
selected summands. In the first line, a horizontal removal
of two boxes from \((D+1,1,0,\ldots,0)\), allowing a
negative last coordinate, has decrements in the first,
second and last coordinates whose sum is two. The second
coordinate can decrease by at most one. This gives exactly
\((0,0,2),(1,0,1),(0,1,1),(2,0,0),(1,1,0)\),
the five terms displayed. All intermediate zero coordinates
stay zero by the interlacing condition. The one-box removals
give the next two lines. The assumptions \(d\geq3\) and
\(D\geq3\) ensure all of these are distinct dominant
weights, including at \(d=3\).
Dual Pieri follows from ordinary Pieri by duality and a
determinant twist; the ordinary rule is Macdonald, I(5.16)--(5.17),
p.~73. Thus no positivity conjecture is used in these identities.

Subtracting the last two lines from the first and adding
\([h_D]\) isolates \([L_\theta]\). Ordinary one-box
Pieri also gives
\[
 S_{(D+1,1)}V=h_{D+1}\otimes V-h_{D+2},\qquad
 S_{(D,1)}V=h_D\otimes V-h_{D+1}
\]
in the same group. Substitution and cancellation yield
\begin{equation}\label{eq:F024character}
 \begin{split}
 [L_\theta]
  ={}&[h_{D+1}\otimes V\otimes\operatorname{Sym}^2V^*]
   -[h_{D+2}\otimes\operatorname{Sym}^2V^*]\\
   &-[h_D\otimes V\otimes V^*]+[h_D].
 \end{split}
\end{equation}
This is an exact character identity; its use below is
justified by exactness of balanced-weight extraction and
finite-group invariants in characteristic zero.

\begin{proof}[The positive coefficient equals one]
Let \(e_1,\ldots,e_d\) be the monomial basis of \(V\),
and let \(\mathbf a=(a,\ldots,a)\). The balanced-weight
subspaces of the four actual tensor modules in
\eqref{eq:F024character} have permutation bases under \(S_d\).
The invariant dimension of such a basis is the number of
its orbits.

For the first term, a basis vector is specified by a marked
index \(j\) in the factor \(V\) and a degree-two dual
multi-index \(\beta\). Its symmetric-power exponent is
forced to be
\[
                           \alpha=\mathbf a+\beta-e_j.
\]
All entries are nonnegative because \(a\geq1\), and
\(|\alpha|=D+1\). Thus every pair \((j,\beta)\)
is allowed. There are exactly four permutation orbits:
\(\beta=2e_j\); \(\beta=2e_k\) with \(k\ne j\);
\(\beta=e_j+e_k\) with \(k\ne j\); and
\(\beta=e_k+e_l\) with \(j,k,l\) all distinct.
The last orbit exists because \(d\geq3\).
Consequently this term contributes four invariant vectors.

For the second term, only \(\beta\) is chosen, and
\(\alpha=\mathbf a+\beta\) is forced. There are two
orbits, according as its two indices are equal or distinct.
For the third term, choose the standard index \(j\) and
the dual index \(k\); the forced exponent is
\(\alpha=\mathbf a+e_k-e_j\), again nonnegative.
There are two orbits, \(j=k\) and \(j\ne k\).
The last term has the single balanced monomial
\(e_1^a\cdots e_d^a\).

The balanced weight is invariant under the permutation
normalizer of the diagonal torus. Taking its subspace is
exact on rational torus modules, and then taking \(S_d\)
invariants is exact over \(\mathbb C\). We may therefore
apply these dimensions to the character identity, obtaining
\[
                        a_\lambda(d[n])=4-2-2+1=1.
\]
This proves \eqref{eq:F024ambient} for all stated parameters.
\end{proof}

\paragraph{The unique copy has an explicit Hessian bracket.}
Put \(s=d-1\geq2\), and let \(F\) be a homogeneous
form of degree \(n\) in \(s\) variables. Define
\[
 H_F(x)=\det\operatorname{Hess}_x F,
           \qquad q=s(n-2),\qquad L=n+q-1,
\]
and the polynomial one-form, written with an auxiliary vector
\(y\),
\[
 \Theta_F(x;y)=qH_F(x)\,dF_x(y)-nF(x)\,dH_{F,x}(y).
                                                               \tag{F-024.5}
\]
It has degree \(s+1=d\) in the coefficients of \(F\),
degree \(L\) in \(x\), and degree one in \(y\).
Euler's formula gives \(\Theta_F(x;x)=0\).
Its coefficient span is consequently the image of
\[
 \det(U)^2\otimes
 \ker\bigl(\operatorname{Sym}^L U\otimes U
                       \longrightarrow\operatorname{Sym}^{L+1}U\bigr)
       =\det(U)^2\otimes S_{(L,1)}U,
                                                               \tag{F-024.6}
\]
where \(U\) is the coordinate-function representation,
and the arrow symmetrizes the two factors. The equality is
ordinary one-box Pieri. More explicitly, the Euler identity
means that the coefficient map kills the fully symmetric
summand; its quotient is identified with the displayed
kernel by that Pieri decomposition. The determinant-square covariance
follows from
\[
 H_{F\circ g}(x)=\det(g)^2H_F(gx),\qquad
 \Theta_{F\circ g}(x;y)=\det(g)^2\Theta_F(gx;gy).
\]
Thus (F-024.6) is an equivariant coefficient map into
\(\operatorname{Sym}^d\operatorname{Sym}^n U\).
The source is irreducible with highest weight
\[
  (L+2,3,2^{s-2})=(d(n-2)+3,3,2^{d-3})=\lambda.
\]
A nonzero such coefficient map has image a complete copy
of this irreducible module.

To check nonvanishing, take \(F=\sum_{i=1}^s x_i^n\),
\(x=(2,1,\ldots,1)\), and \(y=e_1\). Then
\[
 H_F(x)=[n(n-1)]^s2^{n-2}\ne0,
\]
and direct substitution gives
\[
 \Theta_F(x;e_1)
   =\frac{n(n-2)(s-1)(2^n-1)}2 H_F(x)\ne0.
                                                               \tag{F-024.7}
\]
Hence the coefficient span of \(\Theta\) is a nonzero
copy of \(S_\lambda U\). By the ambient multiplicity
calculation it is the unique copy. In particular, to prove
that the corresponding ideal multiplicity is one, or that
the corresponding coordinate-ring multiplicity is one, it
suffices to prove respectively the vanishing or nonvanishing
of this explicit bracket on the \(\mathrm{GL}(U)\)-stable
variety in question.
That final geometric test is separate from the ambient
coefficient computation.

\paragraph{Exact computation and geometric dictionary.}
As a separate finite check, the supplied script
\path{agents/jet_review/check_hessian_bracket_plethysm.py}
expands \(h_d[h_3]\) by its wreath-product cycle index
and evaluates the Schur coefficients by exact
Murnaghan--Nakayama characters. It gives one for
\(s=2,\ldots,12\), including the nine-variable case
\((13,3,2^7)\) at outer degree ten. Its beta-bead
character engine was compared with the independent connected
border-strip engine on all 919 shape/cycle pairs of sizes
zero through eight. These computations supplement the full
arbitrary-parameter proof above.

The torus grading and finite-group action in (F-024.2)
are comodules in the sense of SGA~3, I,
Propositions~4.7.2--4.7.3, pp.~34--36. The Schur and
exterior constructions have the exact-functor dictionary
of SGA~6, VI, Theorem~3.3, pp.~392--393.
The intermediate rational twist is an invertible character
of \(\mathrm{GL}_d\); it is not an effective polynomial
representation on its own. Exactness is used only in
characteristic zero, where the permutation-group invariants
are exact. The Hessian bracket is an actual polynomial
coefficient morphism, rather than a virtual class, and its
vanishing is an ideal-comodule statement. These distinctions
make the passage from the calculated ambient multiplicity
to an actual orbit-closure comparison precise.
